\section{Descent of positivity through the Iitaka fibration}
\label{in:descent}

We now derive a contradiction from Proposition~\ref{in:fibre-big}.
The proof keeps the boundary of actual meridian orders on every
prepared model.  An auxiliary boundary will be used only to study
Iitaka fibres; the final bigness statement will again concern the
actual boundary.

\subsection{A minimum over models and isogenies}

For every prepared model, including those after torus isogenies,
Assumption~\ref{in:input-tori} and Proposition~\ref{in:fibre-big} give
\begin{equation}\label{in:adjoint-lower-bound}
 \kappa(Z,K_Z+\Delta_Z)\ge\ell>0.
\end{equation}
Its hypotheses hold because $Z$ is smooth projective, $g:Z\to A$
is a fibration onto a torus, and $\Delta_Z$ is effective SNC with
coefficients strictly less than one.  Minimize the integer on the
left over \emph{all} these prepared models and all the indicated
isogenies.  The minimum is attained; denote it by $b$ and choose an
attaining model $Z_0$.  Relabel the torus, source, and quotients after
the isogeny producing that model.  Then
\begin{equation}\label{in:adjoint-minimum}
 b=\kappa(Z_0,K_{Z_0}+\Delta_{Z_0})\ge\ell,
\end{equation}
and the lower bound $b$ remains valid on every further prepared model
and after every further torus isogeny.  The source root orbifold
remains special.

Resolve the Iitaka fibration and all relevant supports on a smooth
projective modification $\mu:Z_1\to Z_0$, obtaining
\[
 h:Z_1\longrightarrow W,
 \qquad \dim W=b,
 \qquad g_1:Z_1\longrightarrow A,
\]
where $W$ is smooth projective and both maps are fibrations.  The
model may also resolve a plurisection system realizing the Iitaka
dimension.  Choose a rational SNC boundary $\Theta$ on $Z_1$ such
that
\begin{equation}\label{in:auxiliary-boundary}
 \Theta\ge\Delta_{Z_1},\qquad
 L:=K_{Z_1}+\Theta
 =\mu^*(K_{Z_0}+\Delta_{Z_0})+E,
 \qquad E\ge0\text{ exceptional},
\end{equation}
and all coefficients of $\Theta$ belong to $[0,1)$.

Here are the coefficient choices explicitly.  Keep the coefficients
of the strict transform of $\Delta_{Z_0}$, which equal the
corresponding actual coefficients. In the discrepancy formula
\[
 K_{Z_1}+\mu^{-1}_*\Delta_{Z_0}
 =\mu^*(K_{Z_0}+\Delta_{Z_0})+\sum_F a_F F,
\]
every exceptional $a_F$ is greater than $-1$ because the original
pair is klt.  The actual boundary coefficient $d_F$ is also less
than one.  Choose a rational $t_F<1$ satisfying
$t_F\ge\max(0,d_F,-a_F)$ and put that coefficient on $F$.
Then $a_F+t_F\ge0$, giving \eqref{in:auxiliary-boundary}.

Adding this effective exceptional divisor does not change the
plurisection spaces in divisible degrees.  A section upstairs pushes
to a meromorphic section downstairs with no poles at any prime
divisor, and extends across codimension two; the reverse inclusion is
pullback followed by multiplication by the exceptional section.
Thus $\kappa(Z_1,L)=b$, and $h$ is still its Iitaka fibration.
Moreover $L$ is big on a very general $g_1$-fibre: on such a fibre
\eqref{in:auxiliary-boundary} is a birational pullback of the big
adjoint on the corresponding $Z_0$-fibre, plus an effective divisor.

We give the two restriction arguments needed next.

\begin{lemma}\label{in:iitaka-fibre}
For a very general smooth fibre $H$ of $h$,
\[
 \kappa(H,L|_H)=0,
 \qquad L|_H=K_H+\Theta|_H.
\]
The restricted pair is smooth projective and klt.
\end{lemma}

\begin{proof}
Adjunction and the boundary assertion hold on very general smooth
fibres, using generic smoothness also on the boundary strata.
Some positive multiple of $L$ has a nonzero section which is not
identically zero on a very general fibre, so the restricted Iitaka
dimension is nonnegative.

There are an ample divisor $B$ on $W$, a positive integer $N$, and
an effective divisor $F$ with
\begin{equation}\label{in:ample-remainder}
 NL\sim_{\mathbb Q}h^*B+F,
\end{equation}
after clearing denominators.  Indeed, resolve a complete system of
maximal image dimension.  Its moving part comes from the pullback of
a hyperplane divisor on its image, which gives a big divisor on the
Stein base and on a smooth model of that base.  A sufficiently high
multiple of this big divisor contains an ample divisor.  The fixed
part and the effective remainder give \eqref{in:ample-remainder}.

Suppose that a positive multiple of $L|_H$ had ratios varying along
a very general $H$.  Choose this fibre also in the base-change loci
of all the countably many direct images $h_*\mathcal O_{Z_1}(mL)$
in divisible degrees.  For the relevant degree, twisting that
coherent direct image by a sufficiently high power $\mathcal O_W(kB)$
makes it globally generated.  The resulting sections of
$mL+h^*(kB)$ have ratios varying on general fibres.  Multiplying
them by a section of $k(NL-h^*B)$ from
\eqref{in:ample-remainder} puts them in $(m+kN)L$ and leaves their
ratios unchanged on general fibres.  Combining these sections with a
system having $b$ independent base ratios gives a plurisection map
of dimension at least $b+1$: products place the two systems in a
common degree, retaining both their base ratios and a nonconstant
fibre ratio.  This contradicts $\kappa(Z_1,L)=b$.
\end{proof}

\begin{lemma}\label{in:common-fibres}
Let $g:V\to A_0$ and $h_0:V\to W_0$ be surjective holomorphic maps
of smooth projective varieties with connected general fibres, and let
$M$ be a rational line bundle on $V$ which is big on a very general
smooth $g$-fibre.  Then $M$ is big on every component of a general
fibre of
\[
 (h_0,g):V\longrightarrow J\subset W_0\times A_0.
\]
Equivalently, one may take a very general $h_0$-fibre first and then
a general fibre of its map to its image in $A_0$.
\end{lemma}

\begin{proof}
Take a very general smooth $g$-fibre where base change holds for the
direct images of all divisible multiples of $M$.  Bigness supplies
one degree for which its section map has full generic differential
rank.  On an open of $A_0$ where the corresponding direct image is
locally free, relative evaluation therefore has injective vertical
differential on a dense open $U\subset V$.  The rank condition is
intrinsic and independent of a local frame of that direct image.

Every component of a general fibre of $V\to J$ meets $U$.  To see
this, apply the fibre-dimension theorem to the finitely many
components of the proper closed complement $V\setminus U$.  A
component dominating $J$ has strictly smaller general fibre dimension
than $V\to J$, and a component not dominating $J$ is avoided over
a general point.  The fibres themselves are equidimensional after
shrinking $J$.

At a general smooth point of each such fibre component, its tangent
space is contained in the $g$-vertical tangent space.  The relative
sections restrict to sections of $M$ on the component and still have
injective differential there.  They give a generically finite map,
which proves bigness.  The same open in $J$ meets the slices over
very general points of $W_0$ in dense opens, proving the last
formulation.
\end{proof}

\subsection{The Iitaka fibres are birational to torus translates}

\begin{lemma}\label{in:torus-fibres}
For a very general $h$-fibre $H$, its Albanese morphism
$\beta_H:H\to\operatorname{Alb}(H)$ is birational, and
$g_1|_H$ is the composition of $\beta_H$ with a translated isogeny
onto a subtorus translate in $A$.  Every component of $\Theta|_H$
has image of codimension at least two in that translate.
\end{lemma}

\begin{proof}
Lemma~\ref{in:iitaka-fibre} and
Assumption~\ref{in:input-tori}, applied with the point torus if
necessary, make $\beta_H$ surjective with connected fibres.
The map to $A$ factors through it by a translated homomorphism
$\operatorname{Alb}(H)\to A$.  Quotient $\operatorname{Alb}(H)$
by the connected component of that homomorphism's kernel.  We obtain
a fibration $H\to U_H$ onto a torus, followed by a translated
isogeny from $U_H$ onto the image of $H$ in $A$.

Apply Lemma~\ref{in:common-fibres} to $g_1,h$ and $L$.  It says
that $L|_H$ is big on the components of general fibres of $g_1|_H$,
and hence on a general fibre of $H\to U_H$.  On that smooth fibre
its restriction is the adjoint of the restricted klt pair.
Assumption~\ref{in:input-tori} on $H\to U_H$ now gives
\[
 0=\kappa(H,K_H+\Theta|_H)
 \ge\dim H-\dim U_H.
\]
The fibre dimension is therefore zero.  Since $H$ surjects onto
$\operatorname{Alb}(H)$ and that torus surjects onto $U_H$, all
three dimensions are equal.  Connected general fibres make
$\beta_H$ birational, and the homomorphism from its Albanese torus
to its image in $A$ has finite kernel.

Choose the effective relative canonical divisor for the birational
morphism $\beta_H$; its coefficient is positive on every exceptional
prime, by the vanishing of the Jacobian over a smooth target.  If a
component of $\Theta|_H$ mapped to a divisor $B_H$ in the Albanese
torus, a sufficiently small positive rational multiple of
$\beta_H^*B_H$ would be bounded by $K_H+\Theta|_H$.  At the strict
transform this uses its positive boundary coefficient, and at the
finitely many exceptional primes it uses the positive relative
canonical coefficients.  A nonzero effective divisor on an abelian
variety has positive Iitaka dimension, by the effective-divisor
structure used in Proposition~\ref{in:kernel}.  This contradicts
Lemma~\ref{in:iitaka-fibre}.  Thus the boundary is exceptional over
the Albanese torus.  The remaining map is finite, so its images also
have codimension at least two in the image translate.
\end{proof}

The subtorus occurring here is fixed as $H$ varies over very general
fibres.  More precisely, over a connected smooth-fibration open of
$W$, the image
\[
 M_0=\operatorname{im}\bigl(H_1(H,\mathbb Z)\longrightarrow
 H_1(A,\mathbb Z)\bigr)
\]
is constant by smooth proper transport and homotopy of the maps to
$A$.  Its real span determines a fixed subtorus $T_0\subseteq A$.
By Lemma~\ref{in:torus-fibres}, $M_0$ has finite index in
$\Lambda_{T_0}=\pi_1(T_0)$ and records the isogeny on a very general
fibre.  Choose $k>0$ with
\[
 k\Lambda_{T_0}\subseteq M_0.
\]
Make the multiplication isogeny $A^+\to A$ given by $[k]$, and let
$T^+\subset A^+$ be the identity component over $T_0$.  The
subtorus lattice is primitive in the ambient lattice, so the image
of $\pi_1(T^+)$ is $k\Lambda_{T_0}$.

Each connected component of the pullback of a very general $H$
maps birationally onto a translate of $T^+$.  Indeed, the covering
of that translate factors through the isogeny from
$\operatorname{Alb}(H)$ because of the lattice inclusion just
chosen.  Its fibre product with that isogeny is a disjoint union of
degree-one covers of the lifted translate.  The birational Albanese
morphism pulls back birationally to each of them.  There are no extra
components supported over its exceptional locus, since the map to
$H$ is finite \'etale and every component meets the dense open where
the birational morphism is an isomorphism.  The pulled-back boundary
on these fibres still has image of codimension at least two in the
lifted translate.

\subsection{A torus-bundle model and its actual boundary}

Put
\[
 \widetilde Z_1=Z_1\times_A A^+,
 \qquad X^+=X\times_A A^+.
\]
Both spaces are connected smooth projective by
\eqref{in:isogeny-groups}.  The root orbifold on $X^+$ is special.
The actual boundary on $\widetilde Z_1$ is the pullback of
$\Delta_{Z_1}$ and is dominated by the pullback of $\Theta$.

Let $W'$ be the Stein base of $\widetilde Z_1\to W$.  It is finite
over $W$.  The map to $A^+/T^+$ is constant on very general connected
Stein fibres, so it factors rationally through $W'$.  Resolve to a
smooth projective $W_1$ with a morphism to $A^+/T^+$, and form
\begin{equation}\label{in:torus-bundle}
 Z_*=W_1\times_{A^+/T^+}A^+,
 \qquad p:Z_*\longrightarrow W_1.
\end{equation}
The natural rational comparison from $\widetilde Z_1$ is birational:
on a very general connected Stein fibre it is the degree-one map
onto the lifted torus translate just established.  The map $p$ is
smooth with connected irreducible torus fibres.  Its relative canonical
bundle is trivial, since it is the pullback of a quotient map of
tori and the invariant relative volume form is nowhere zero.  Hence
\begin{equation}\label{in:canonical-bundle}
 K_{Z_*}=p^*K_{W_1}
\end{equation}
with compatible canonical divisors.  The construction works after
any further smooth projective modification of $W_1$.

Let $R^+=G^+/B_0\le R$ be the quotient for the pulled-back source.
Transfer its open monodromy to $Z_*$ and read the actual boundary
$\Delta_*$.  The same quotient is used on common open subsets of
all comparison models.

\begin{lemma}\label{in:vertical-boundary}
The support of $\Delta_*$ is vertical over $W_1$.  Consequently
there is an effective rational divisor $\Delta_{W_1}$, with finite
integral root coefficients, such that
\begin{equation}\label{in:boundary-pullback}
 \Delta_*=p^*\Delta_{W_1}.
\end{equation}
The coefficients on corresponding primes record the same actual
orders.  These assertions also hold after any further smooth
projective modification of $W_1$, with the orders read anew.
\end{lemma}

\begin{proof}
Compare $Z_*$ and $\widetilde Z_1$ on a common smooth projective
resolution, with compatible maps to $A^+$ and to $W'$.  Suppose a
prime of nontrivial order on $Z_*$ were horizontal over $W_1$.
Its strict transform maps, over a very general parameter, onto a
divisor in the corresponding torus translate in $A^+$.  If the
strict transform is exceptional over $\widetilde Z_1$, its centre
has codimension at least two there.  On very general fibres over
$W'$ that centre has codimension at least two or is absent, by the
fibre-dimension theorem.  Its image in the lifted translate cannot
be the indicated divisor.

If it is not exceptional, it corresponds to a divisor of nontrivial
actual order on $\widetilde Z_1$.  This divisor belongs to the
pullback of $\Theta$, because that pullback dominates the actual
boundary.  By Lemma~\ref{in:torus-fibres} and the \'etale pullback
description, its restriction to very general connected Stein fibres
has image of codimension at least two in the lifted torus translate.
Again it cannot supply a divisor there.  This contradiction excludes
horizontal support.  When $T^+$ has dimension zero, horizontal
prime divisors do not exist and the conclusion is immediate.

A vertical prime of the smooth map $p$ must dominate a prime of
$W_1$: the fibre dimension bound gives codimension at most one for
its image.  Over that prime the full inverse image is irreducible,
since its general torus fibre is geometrically irreducible, and it
is reduced in codimension one by smoothness.  Our vertical prime is
therefore exactly that inverse image, with pullback multiplicity
one.  Assign its actual coefficient to the corresponding prime on
$W_1$ to obtain \eqref{in:boundary-pullback}.  The same comparison
argument uses very general fibres over $W'$ after any further
modification of $W_1$, and proves the last assertion.
\end{proof}

\begin{proposition}\label{in:base-big}
For the actual boundary in \eqref{in:boundary-pullback},
\[
 K_{W_1}+\Delta_{W_1}\text{ is big}.
\]
This remains true for the version constructed on any further smooth
projective modification of $W_1$.
\end{proposition}

\begin{proof}
Dominate $Z_*$ by a prepared smooth projective model over $A^+$.
Its adjoint, defined by the actual boundary, has Iitaka dimension
at least $b$ by the minimum
\eqref{in:adjoint-minimum}, which included all further isogenies.
For every sufficiently divisible degree, its plurisections inject by
birational pushdown into the plurisections of $K_{Z_*}+\Delta_*$.
At every prime on $Z_*$, its strict transform has the same actual
meridian order.  Thus a meromorphic pluricanonical form pushed
downward has the prescribed pole bound at every prime.  It extends
as a section of the indicated line bundle across codimension two.
This proves the injection and preserves all section ratios.

Equations \eqref{in:canonical-bundle} and
\eqref{in:boundary-pullback}, together with
$p_*\mathcal O_{Z_*}=\mathcal O_{W_1}$, now give
\[
 b\le\kappa(Z_*,K_{Z_*}+\Delta_*)
   =\kappa(W_1,K_{W_1}+\Delta_{W_1}).
\]
Since $\dim W_1=\dim W=b$, the adjoint on $W_1$ is big.  For a
further modification of $W_1$, use its corresponding torus-bundle
model, Lemma~\ref{in:vertical-boundary}, and a new prepared model
dominating it.  The same minimum and pushdown argument applies.
No birational invariance of this actual adjoint is asserted or needed.
\end{proof}

\subsection{Preparing the base and constructing the differential line}

We finish by returning to the fixed smooth root source on $X^+$.
We shall only test primes on that source, and prepare the base to
ensure that the unresolved higher-codimension images do not contribute
additional prime-divisor conditions.

\begin{lemma}\label{in:neatness}
For a rational fibration $u_0:P\dashrightarrow W_0$ of smooth
projective varieties, one can modify the smooth projective base and
resolve the source to obtain
\[
 P\xleftarrow{\tau}\widehat P\overset{u}{\longrightarrow}W_1
\]
such that every prime divisor of $\widehat P$ whose image under $u$
has codimension at least two is exceptional over $P$.
\end{lemma}

\begin{proof}
First resolve the graph to a projective morphism over $W_0$.  Apply
projective flattening by a modification of the integral base, taking
the main transform of that graph; it is flat of the expected relative
dimension over the modified base.  Resolve the base.  The main
transform after this further base change is contained in the pullback
of the flat family, so its fibres have dimension at most the same
expected relative dimension.  It still maps properly and birationally
to $P$.

Resolve this main transform to $\widehat P$.  A prime mapping to
base codimension at least two cannot have a divisorial image in the
main transform, by the fibre dimension bound.  Its image there has
codimension at least two, and its image in $P$ consequently has
dimension at most $\dim P-2$.  It is exceptional over $P$, as
required.  All steps are projective in characteristic zero; the
flattening and base-change form used here is standard
\cite{Stacks,Hironaka}.
\end{proof}

\begin{proof}[Proof of Theorem~\ref{in:main}]
Apply Lemma~\ref{in:neatness} to the induced rational fibration
$X^+\dashrightarrow W_1$, allowing a further modification of $W_1$.
Proposition~\ref{in:base-big} is valid for its new actual boundary.
Write
\[
 X^+\xleftarrow{\tau}\widehat X\overset{u}{\longrightarrow}W_1.
\]
The maps to $W_1$ and to $A^+$ give a holomorphic map
$\widehat X\to Z_*$ by the fibre product
\eqref{in:torus-bundle}: their composites to $A^+/T^+$ agree on a
dense open and hence everywhere.  It is compatible with the open
monodromy to $R^+$.

Let $E$ be the strict transform of a prime of $X^+$, let $r_E=2$
when that prime belongs to its root boundary and $r_E=1$ otherwise.
If $E$ dominates a boundary prime $B$ of $W_1$, write $m_B$ for
its order and $e=\operatorname{ord}_E(u^*B)$.  Lemma~\ref{in:local-divisibility},
applied to the smooth torus bundle and
\eqref{in:boundary-pullback}, gives
\begin{equation}\label{in:final-divisibility}
 m_B\mid e r_E.
\end{equation}
This also covers a source divisor whose image in $Z_*$ is a proper
subset of $p^{-1}(B)$.  Around that image over a general point of
$B$, filling the other deleted loci kills only meridians of trivial
image.  The remaining local group is the complement group of the
one smooth hypersurface $p^{-1}(B)$.  Its generator has order $m_B$;
a transverse source disc winds $e$ times and its boundary has order
dividing $r_E$ in $G^+$.  This explains precisely why the weighted
divisibility applies at these possibly exceptional intermediate
images.

Take the saturated line $\mathcal L$ in
$\Omega^b_{\mathcal X^+}$ determined generically by the top
differential of the rational map to $W_1$.  It is an orbifold line
bundle by the saturation argument in Section~\ref{in:foundations}.
We claim that, for all sufficiently divisible $j$, pullback gives
\begin{equation}\label{in:final-sections}
 H^0\bigl(W_1,\mathcal O_{W_1}(j(K_{W_1}+\Delta_{W_1}))\bigr)
 \longrightarrow H^0(\mathcal X^+,\mathcal L^{\otimes j}).
\end{equation}
It suffices to check regularity at every prime in the smooth root
charts.  At the general point corresponding to $E$, the modification
$\tau$ is an isomorphism.  If $E$ dominates a boundary prime as
above, a local equation of that prime has order $e r_E$ in the
source root coordinate.  The pulled-back local top differential has
common vanishing at least $e r_E-1$.  By
\eqref{in:final-divisibility},
\[
 e r_E-1\ \ge\ e r_E(1-m_B^{-1}),
\]
which cancels the pole permitted by the coefficient of $B$ in each
tensor factor.  Saturation makes this common divisorial factor the
relevant test for a section of the line power.

At primes dominating $W_1$, or mapping generically outside its
boundary, there is no boundary pole.  A strict source prime cannot
map into base codimension at least two by Lemma~\ref{in:neatness}.
These exhaust the primes of the original $X^+$.  Extension across
codimension two on its smooth root charts proves
\eqref{in:final-sections}.  Thus no specialness assertion is needed
on the resolved source $\widehat X$.

Pullback along a dominant map preserves ratios.  Bigness from
Proposition~\ref{in:base-big} therefore implies
$\kappa(\mathcal X^+,\mathcal L)\ge b$.  The root differential
bound of Lemma~\ref{in:root-bound} gives the reverse inequality.
Since $1\le\ell\le b\le\dim X^+=4$, this is a forbidden
differential line on the special orbifold $\mathcal X^+$.

The contradiction excludes the infinite-kernel case in
Proposition~\ref{in:group-reduction}.  The reductions in
Section~\ref{in:finite-cover-cases} handle a trivial meridian and a
meridian detected in a finite quotient by applications of
Assumption~\ref{in:input-manifold} to
the appropriate special smooth projective manifolds.  A nontrivial
meridian invisible in all finite quotients would give exactly the
infinite kernel excluded above.  Hence the original entire orbifold
group $G$ has an abelian subgroup of finite index.
\end{proof}
